\section{The invariant Hardy-gauge contour matrix}
\label{sec:common-matrix}

\subsection{The analytic Hardy gauge}

Fix once and for all
\[
c:=2.
\]
For horizontal heights $T_1<T_2$ avoiding the zeros relevant to the contour,
let
\begin{equation}
\mathcal R_T^{\rm sym}
:=\{\sigma+it:-1\le\sigma\le2,\ T_1\le t\le T_2\}.
\label{eq:symmetric-contour-rectangle}
\end{equation}
The existence of admissible $T_1\in[T,T+1]$ and
$T_2\in[2T,2T+1]$ is proved in
Proposition~\ref{prop:z1-zeta-decrement}. On a simply connected neighbourhood
of this symmetric high rectangle choose a holomorphic branch
\[
q(s):=\chi(1-s)^{1/2}.
\]
Fix its sign at one point of the critical line so that the resulting Hardy
gauge is real there, and put
\[
\mathscr Z(s):=q(s)\zeta(s).
\]
This normalization determines the branch on the rectangle. Indeed the
holomorphic function
\[
q^\dagger(s):=\overline{q(1-\bar s)}
\]
satisfies $(q^\dagger)^2=\chi(s)=q(s)^{-2}$; the chosen sign gives
$q^\dagger=q^{-1}$ at the base point and hence everywhere. Using the zeta
functional equation therefore yields the exact reflection law
\begin{equation}
\boxed{\mathscr Z(1-\bar s)=\overline{\mathscr Z(s)}.}
\label{eq:hardy-gauge-reflection}
\end{equation}
In particular $F(z)=\mathscr Z(1/2+iz)$ satisfies
$F(\bar z)=\overline{F(z)}$ and is real on the real axis. The factor $q$ is
analytic and nonvanishing, so $\mathscr Z$ and $\zeta$ have exactly the same
zeros with the same multiplicities. Define
\[
f(s):=-\frac12\frac{\chi'}{\chi}(s),
\qquad
Z_1(s):=\zeta'(s)+f(s)\zeta(s).
\]
Since $q'/q=f$, differentiation gives the exact identity
\begin{equation}
\boxed{\mathscr Z'(s)=q(s)Z_1(s).}
\label{eq:gauge-derivative}
\end{equation}
Thus the stationary zeros of the Hardy-gauge detector are precisely the zeros
of the same $Z_1$ which appears on the arithmetic side.

\subsection{Curvature representation of the contour form}

Let $a$ be the physical shift and put
\[
p:=La.
\]
The curvature source is generated before any right-edge expansion from one
shifted Hardy product. Algebraically,
\begin{equation}
\mathscr Z(s)\mathscr Z''(s)-\mathscr Z'(s)^2
=\frac12\left.\partial_a^2\bigl(\mathscr Z(s+a)\mathscr Z(s-a)\bigr)\right|_{a=0}.
\label{eq:curvature-product}
\end{equation}
Passing from $a$ to $p$ introduces the positive scalar factor $L^2$. We divide
the raw matrix by this factor once and for all. Since the rescaling is
positive, it does not change inertia.

\begin{proposition}[Exact normalized same-form identity]
\label{prop:same-form-identity}
Let
\[
C_{\mathscr Z}(s):=\mathscr Z(s)\mathscr Z''(s)-\mathscr Z'(s)^2,
\qquad
C_{\mathscr Z}^{(p)}(s):=L^{-2}C_{\mathscr Z}(s).
\]
Equivalently,
\[
C_{\mathscr Z}^{(p)}(s)
=\frac12\left.\partial_p^2
\bigl(\mathscr Z(s+p/L)\mathscr Z(s-p/L)\bigr)\right|_{p=0}.
\]
Let $\{\psi_\nu\}_{\nu\in\mathcal I_T}$ be the smooth entire packet family
specified below and write
\[
B_v(z):=\sum_{\nu\in\mathcal I_T}v_\nu\psi_\nu(z),
\qquad B_v^\#(z):=\overline{B_v(\bar z)}.
\]
With $z(s):=-i(s-\tfrac12)$ define the holomorphic polarization
\begin{equation}
\boxed{\mathcal B_{v,w}(s):=B_w(z(s))B_v^\#(z(s)).}
\label{eq:packet-polarization}
\end{equation}
Then the normalized invariant contour form is exactly
\begin{equation}
\boxed{
\mathcal A_T(v,w)
=\frac1{2\pi i}\oint_{\partial\mathcal R_T^{\rm sym}}
\frac{\mathscr Z''}{\mathscr Z'}(s)
C_{\mathscr Z}^{(p)}(s)\mathcal B_{v,w}(s)\,ds
=\frac1{2\pi i}\oint_{\partial\mathcal R_T^{\rm sym}}
\frac{Z_1'}{Z_1}(s)
C_{\mathscr Z}^{(p)}(s)\mathcal B_{v,w}(s)\,ds.
}
\label{eq:same-form-identity}
\end{equation}
After straightening the critical line by $s=1/2+iz$, the same form is exactly
\eqref{eq:zero-kernel}. Thus the zero-side stationary residue detector and the
arithmetic $Z_1'/Z_1$ source are literally the same Hermitian contour matrix.
\end{proposition}

\begin{proof}
Since \eqref{eq:gauge-derivative} gives
\[
\frac{\mathscr Z''}{\mathscr Z'}
=f+\frac{Z_1'}{Z_1},
\]
and $f$, $C_{\mathscr Z}^{(p)}$ and $\mathcal B_{v,w}$ are holomorphic on the
high rectangle, the integral of
$fC_{\mathscr Z}^{(p)}\mathcal B_{v,w}$ around the closed contour is zero. This proves
the second equality in \eqref{eq:same-form-identity}.

For the first equality, expand
\[
\frac{\mathscr Z''}{\mathscr Z'}C_{\mathscr Z}^{(p)}
=L^{-2}\left(
\frac{\mathscr Z\,\mathscr Z''^2}{\mathscr Z'}
-\mathscr Z'\mathscr Z''\right).
\]
The second term, after multiplication by the holomorphic packet polarization,
is holomorphic and hence has zero closed-contour integral. Put
$F(z)=\mathscr Z(1/2+iz)$. Then
\[
\mathscr Z'=-iF',\qquad \mathscr Z''=-F'',\qquad ds=i\,dz,
\]
and therefore the remaining term transforms to
\[
\frac1{2\pi i}\oint
L^{-2}\left(-F(z)\frac{F''(z)^2}{F'(z)}\right)
 B_w(z)B_v^\#(z)\,dz,
\]
which is exactly \eqref{eq:zero-kernel}. Finally, \eqref{eq:curvature-product}
shows that this same form is generated from one shifted Hardy product before
any contour movement or stationary reduction.
\end{proof}

With the packet notation above, the normalized Hermitian contour form is
\begin{equation}
\mathcal A_T(v,w)
:=\frac1{2\pi i}\int_{\partial\mathcal R_T^{\rm sym}}
\mathcal K(s;B_v,\overline{B_w})\,ds,
\label{eq:contour-form}
\end{equation}
where $\mathcal K$ is the exact Hardy-gauge curvature integrand obtained from
\eqref{eq:curvature-product}, before Dirichlet expansion, stationary phase,
Perron inversion, shell analysis, or arithmetic contraction. Define the finite matrix $A_T$ by
\[
\langle A_Tv,w\rangle=\mathcal A_T(v,w).
\]
The reflection law \eqref{eq:hardy-gauge-reflection} together with
\eqref{eq:packet-polarization} gives
\[
\mathcal A_T(w,v)=\overline{\mathcal A_T(v,w)},
\]
so $A_T$ is Hermitian. It is the invariant Hermitian object throughout the
argument. The positive form $\Gfr_T$ introduced later is not a second invariant
object; it is the fixed reference metric used for relative operator estimates.

\begin{lemma}[Invariance under standard reductions]
\label{lem:allowed}
Every later change of representation used below belongs to one of the
following four classes.
\begin{enumerate}[label=(\roman*)]
\item An exact congruence $A\mapsto U^*AU$ with $U$ invertible; inertia is
preserved exactly.
\item Restriction to a subspace of codimension $d$; at most $d$ negative
directions are lost.
\item Addition of a matrix of rank $r$; negative inertia changes by at most
$r$.
\item Addition of a Hermitian remainder $R$ relative to a positive reference
Gram $G$, measured by
\[
\|G^{-1/2}RG^{-1/2}\|_{S_2}^2.
\]
\end{enumerate}
\end{lemma}

\begin{proof}
The first assertion is Sylvester's law of inertia. The second is the elementary
dimension inequality for the intersection of a negative subspace with a
codimension-$d$ subspace. The third follows by applying the same argument to
$\ker E$, whose codimension is at most $r$. The fourth is quantified in
Lemma~\ref{lem:minmax} below.
\end{proof}
